\section{今天你的 Harness 是什么？}

\subsection{渐进式构建 Harness 的起点}

OpenAI 团队花了 5 个月来构建他们的 harness------这不是一蹴而就的事情。但 Birgitta 鼓励每个团队反思自己当前的 harness 是什么，并提出了一系列引导性问题：

\begin{importantbox}{Harness 自检清单}
以下问题可以帮助你评估当前团队的 harness 成熟度：
\begin{enumerate}[leftmargin=1.5em]
    \item \textbf{你有 pre-commit hook 吗？}里面包含了什么检查？
    \item \textbf{你有自定义 linter 的想法吗？}有哪些项目特定的规则想要自动化？
    \item \textbf{你想在代码库上施加哪些架构约束？}模块边界是否清晰？
    \item \textbf{你尝试过结构性测试框架吗？}比如 ArchUnit 这样验证架构规则的工具？
\end{enumerate}
即使你还没有使用 AI agent 进行开发，这些问题本身就代表了良好的工程实践。
\end{importantbox}

\subsection{从人类实践到 AI 实践的连续性}

这个自检清单揭示了一个重要的事实：harness engineering 的许多组件------linter、结构性测试、架构约束------本质上就是\textbf{良好的软件工程实践}。AI agent 的出现并没有发明这些概念，而是赋予了它们新的紧迫性和重要性。

在传统开发中，我们可以依靠经验丰富的开发者通过代码审查来执行隐式的架构约定。但当 AI agent 成为主要的代码生产者时，这些约定必须被\textbf{显式化}和\textbf{自动化}------因为 agent 无法理解``这里我们团队一般不这么做''这样的隐性知识。

\begin{knowledgebox}{ArchUnit 和结构性测试}
ArchUnit 是一个用于检查 Java/Kotlin 代码架构规则的测试框架。它允许你编写类似以下的规则：``所有 Service 类不应依赖 Controller 类''、``所有 Repository 接口必须位于 persistence 包中''等。类似的工具在其他语言中也有对应：
\begin{itemize}[leftmargin=1.5em]
    \item Java/Kotlin: ArchUnit
    \item .NET: NetArchTest
    \item Python: import-linter
    \item TypeScript: eslint-plugin-boundaries
\end{itemize}
这些工具是构建 harness 架构约束层的重要基础。
\end{knowledgebox}

\subsection{本章小结}

构建 harness 是一个渐进的过程，可以从评估现有的 pre-commit hook、linter 和架构约束开始。harness 的核心组件本质上是良好工程实践的显式化和自动化。即使尚未使用 AI agent，建立这些实践也能提升团队的工程质量。

\newpage

